\documentclass[10pt,a4paper]{amsart}
\usepackage[margin=19mm]{geometry}
\usepackage{amsmath,amssymb,mathtools}
\usepackage[scr=rsfs,cal=euler]{mathalfa}
\usepackage{xfrac}
\usepackage[unicode,hidelinks,bookmarks=false]{hyperref}
\newcommand{\ie}{i.e.\ }
\newcommand{\cf}{cf.\ }
\newcommand{\resp}{respectively\ }
\newcommand{\Subsection}{\subsection}
\providecommand{\nobreakdash}{\nobreak\hbox{-}}
\setlength{\emergencystretch}{2em}
\multlinegap0pt
\usepackage{bm}
\usepackage{bbm}
\usepackage{dsfont}
\usepackage{xspace}
\usepackage{cancel}
\usepackage{mathtools}
\usepackage{amsthm}
\makeatletter
\@ifundefined{theorem}{\newtheorem{theorem}{Theorem}}{}
\makeatother
\newcommand{\Rmn}{\mathbb{R}^{m\times n}}
\newcommand{\R}{\mathbb{R}}
\newcommand{\be}{\begin{eqnarray}}
\newcommand{\ee}{\end{eqnarray}}
\title[Incremental Column Subset Selection via Conditional Determin]{Incremental Column Subset Selection via Conditional Determinantal Point Processes}
\author{Laura Grigori, Zhipeng Xue}
\date{Source: arXiv:2609.20556v1 (CC BY 4.0)}
\begin{document}
\maketitle

\noindent\emph{Source and licence.} Incremental Column Subset Selection via Conditional Determinantal Point Processes. Version arXiv:2609.20556v1 (CC BY 4.0), distributed under the Creative Commons Attribution 4.0 International licence, as recorded in the arXiv licence metadata for this exact version and shown on the abstract page https://arxiv.org/abs/2609.20556v1. Full rights notice: licenses/arxiv-2609.20556-CC-BY-4.0.txt.\par\smallskip

\noindent\textit{Editorial scope.} Counted unit: the Theorem whose source label is \texttt{thm:conditional-sharpness} in the article (source0.tex lines 837--844, source label \texttt{thm:conditional-sharpness}), delivered together with the article's own complete original proof of it (source0.tex lines 846--1058, one \texttt{proof} environment of 213 lines). No mathematics is written, repaired or re-derived: statement and proof are the article's own text line for line. Required context retained uncounted: eq:conditional\_dpp\_general (lines 826--834); thm:conditional\_dpp\_general (lines 826--834). Every definition, display and label of those lines is retained unchanged. Omitted from this package and not redistributed: 1--825, 835--836, 845--845, 1059--1698. Nothing inside a retained excerpt is omitted or altered: every retained body is delivered exactly as the article's lines are written, so no omission receipt of any kind accompanies this package. Citations: the retained excerpts carry no citation command at all, so no bibliography entry of the article is transcribed and no reference is redistributed. Reference adaptation: none was needed and none was made; every reference inside the delivered unit resolves inside it, so no unresolved reference or citation is produced. Printed slips kept literal: no printed word, symbol, hypothesis or number of the counted statement or of its proof was altered. Layout and class: the article's own class is \texttt{article} with packages including babel, comment, geometry, cite, subcaption, stmaryrd, algorithm, algpseudocode, amsmath, graphicx, listings, rotating, xcolor, inputenc; the wrapper is the lane's pinned \texttt{amsart} editorial wrapper at 10pt on A4, which restates the theorem-like environments the retained text needs and adds the article's own macro definitions that the retained text needs (\texttt{\textbackslash R}, \texttt{\textbackslash Rmn}, \texttt{\textbackslash be}, \texttt{\textbackslash ee}), with extra wrapper packages (bm, bbm, dsfont, xspace, cancel, mathtools). The delivered numbering is the unit's own. Assets: no retained span contains a figure environment, a graphics-inclusion command, a PDF-page inclusion, a table or a TikZ picture; no image file is copied into this package, the package declares no asset dependency and no image file is redistributed with it. Licence: version arXiv:2609.20556v1 of this article is distributed under the Creative Commons Attribution 4.0 International licence, as recorded in the arXiv licence metadata for this exact version and shown on the abstract page https://arxiv.org/abs/2609.20556v1. A full rights notice is delivered at licenses/arxiv-2609.20556-CC-BY-4.0.txt. Characters: every retained excerpt is pure ASCII, so the delivered engine, XeLaTeX with its default Latin Modern text font, reports no missing character.
\begin{theorem}
\label{thm:conditional_dpp_general}
    Let $A\in \Rmn$, $V \in \R^{n\times (k+p)} $ with $V^TV = I$ and $K = VV^T$. Let $S\subset \{1,2,\dots,n\}$ be an index subset with $|S| = k$ and $V(S,:)$ having  full row rank. If we sample $T\sim \mathrm{DPP}(\widetilde K)$ defined above and then let $U = S\cup T$, we have
    \be
    \label{eq:conditional_dpp_general}
    \mathbb{E}_T [ \|A - A(:,U)V(U,:)^{-T}V^T\|_F^2\mid S ] \leq (1 + p) \|\widetilde{A} -  \widetilde{A}(:,S)(V^T(:,S))^{\dagger} V^T\|_F^2,
    \ee
    where $\widetilde A = A - AVV^T$.
\end{theorem}

\begin{theorem}
\label{thm:conditional-sharpness}
For any integers $k\geq 1$ and $p\geq 1$, there exist a matrix
$A\in\mathbb{R}^{(k+p+1)\times(k+p+1)}$, an orthonormal matrix
$V\in\mathbb{R}^{(k+p+1)\times(k+p)}$, and a prescribed initial
subset $S$ with $|S|=k$ such that the equality of~\eqref{eq:conditional_dpp_general} in
Theorem~\ref{thm:conditional_dpp_general} can be attained exactly.
\end{theorem}

\begin{proof}
Let $Q\in\mathbb{R}^{p\times(p+1)}$ satisfy
\be
QQ^\top=I_p,
\qquad
Q\mathbf{1}_{p+1}=0,
\ee
and define
\be
B=
\begin{bmatrix}
Q\\[1mm]
\dfrac{1}{2\sqrt{p+1}}\mathbf{1}_{p+1}^\top
\end{bmatrix},
\qquad
A=
\begin{bmatrix}
I_k & 0\\
0 & B
\end{bmatrix},
\qquad
V=
\begin{bmatrix}
I_k & 0\\
0 & Q^\top
\end{bmatrix}.
\ee
Let $S=\{1,\ldots,k\},
\,
K=VV^\top$,
and we sample
$T\sim\operatorname{DPP}(\widetilde K)$ as defined above.

Since $Q\mathbf{1}_{p+1}=0$ and $QQ^\top=I_p$, we have
\be
Q^\top Q
=
I_{p+1}
-
\frac{1}{p+1}
\mathbf{1}_{p+1}\mathbf{1}_{p+1}^\top.
\ee
Moreover,
\[
BB^\top
=
\begin{bmatrix}
I_p & 0\\
0 & \dfrac14
\end{bmatrix}.
\]
Hence the singular values of $A$ are $1$ with multiplicity $k+p$ and $1/2$. Therefore,
\[
\left\|A-A_{k+p}\right\|_F^2=\frac14.
\]
Moreover,
\be
\widetilde A
:=
A-AVV^\top
=
\begin{bmatrix}
0 & 0\\
0 &
\begin{bmatrix}
0\\
\dfrac{1}{2\sqrt{p+1}}\mathbf{1}_{p+1}^\top
\end{bmatrix}
\end{bmatrix},
\ee
so that
\be
\|\widetilde A\|_F^2
=
\left\|
A-A_{k+p}
\right\|_F^2= \frac14.
\ee
Since $S=\{1,\ldots,k\}$, we also have
$\widetilde A(:,S)=0$,
which gives
\be
\left\|
\widetilde A
-
\widetilde A(:,S)
\bigl(V^\top(:,S)\bigr)^\dagger
V^\top
\right\|_F^2
=
\left\|
A-A_{k+p}
\right\|_F^2.
\ee

We can simplify $K$ and $\widetilde K$ as
\be
K=
\begin{bmatrix}
I_k & 0\\
0 & Q^\top Q
\end{bmatrix},
\quad
\widetilde K
=
\begin{bmatrix}
0 & 0\\
0 & Q^\top Q
\end{bmatrix}.
\ee
Thus every choice of $T$ contains $p$ indices from the last
$p+1$ indices. Let $\widehat T\subset\{1,\ldots,p+1\}$ denote the
corresponding local indices and set
$
Q_{\widehat T}=Q(:,\widehat T).
$
Since $\operatorname{Null}(Q)=\operatorname{span}
\{\mathbf{1}_{p+1}\}$, every $p$ columns of $Q$ are linearly
independent, and hence $Q_{\widehat T}$ is nonsingular.

For $U=S\cup T$, the block structure gives
\be
A(:,U)V(U,:)^{-\top}V^\top
=
\begin{bmatrix}
I_k & 0\\
0 &
B(:,\widehat T)Q_{\widehat T}^{-1}Q
\end{bmatrix}.
\ee
Writing
\be
B(:,\widehat T)
=
\begin{bmatrix}
Q_{\widehat T}\\[1mm]
\dfrac{1}{2\sqrt{p+1}}\mathbf{1}_p^\top
\end{bmatrix},
\ee
we obtain
\be
B-
B(:,\widehat T)Q_{\widehat T}^{-1}Q
=
\begin{bmatrix}
0\\[1mm]
\dfrac{1}{2\sqrt{p+1}}
\left(
\mathbf{1}_{p+1}^\top
-
\mathbf{1}_p^\top Q_{\widehat T}^{-1}Q
\right)
\end{bmatrix}.
\ee

Let $j$ be the unique index not contained in $\widehat T$.
Since $Q\mathbf{1}_{p+1}=0$,
\be
Q(:,j)
=
-Q_{\widehat T}\mathbf{1}_p,
\ee
and therefore
\be
Q_{\widehat T}^{-1}Q(:,j)
=
-\mathbf{1}_p.
\ee
For every selected column, $Q_{\widehat T}^{-1}Q(:,i)$ is a
column of the identity matrix. Hence
\be
\mathbf{1}_{p+1}^\top
-
\mathbf{1}_p^\top Q_{\widehat T}^{-1}Q
\ee
vanishes on all indices in $\widehat T$ and equals $p+1$ at the
single omitted index $j$. Consequently,
\be
\begin{aligned}
\left\|
A-A(:,U)V(U,:)^{-\top}V^\top
\right\|_F^2
&=
\frac{1}{4(p+1)}(p+1)^2\\
&=
\frac{p+1}{4}.
\end{aligned}
\ee
This holds every choice of $T$, and therefore
\be
\mathbb{E}_T
\left[
\left\|
A-A(:,U)V(U,:)^{-\top}V^\top
\right\|_F^2
\,\middle|\,S
\right]
=
\frac{p+1}{4}.
\ee
Since $\|A-A_{k+p}\|_F^2=1/4$, we conclude that
\be
\mathbb{E}_T
\left[
\left\|
A-A(:,U)V(U,:)^{-\top}V^\top
\right\|_F^2
\,\middle|\,S
\right]
=
(p+1)\left\|A-A_{k+p}\right\|_F^2.
\ee
\end{proof}

\end{document}
